% Generated draft: requires mathematical and source review.
\documentclass[11pt]{article}
\usepackage[margin=25mm]{geometry}
\usepackage{fontspec}
\setmainfont{Latin Modern Roman}
\usepackage{amsmath,amssymb,mathtools,unicode-math}
\setmathfont{Latin Modern Math}
\usepackage{enumitem,needspace,xurl,hyperref}
\setlength{\parindent}{0pt}
\setlength{\parskip}{5pt}
\setlength{\emergencystretch}{3em}
\newcommand{\N}{\mathbb{N}}
\newcommand{\Z}{\mathbb{Z}}
\newcommand{\Q}{\mathbb{Q}}
\newcommand{\R}{\mathbb{R}}
\newcommand{\C}{\mathbb{C}}
\newcommand{\F}{\mathbb{F}}
\newcommand{\A}{\mathbb{A}}
\newcommand{\PP}{\mathbb{P}}
\newcommand{\E}{\mathbb{E}}
\newcommand{\eps}{\varepsilon}
\newcommand{\dd}{\,\mathrm{d}}
\begin{document}
\section*{UnsolvedMath ID: 30005261}
\textbf{Source problem code:} catalogue entry
\subsection{Definitions and mathematical statement}
A topological space \(X\) is contractible if there exist \(x_0\in X\) and a continuous map \(H:X\times[0,1]\to X\) such that
\[
H(x,0)=x,\qquad H(x,1)=x_0\qquad\text{for every }x\in X.
\]
This definition explains the term; the supplied record does not identify the space or topology to which the hypothesis in the cited theorem applies.

\textbf{Source-relative question.} In Theorem 3 of the cited report contribution, can the contractibility hypothesis be omitted while retaining every other hypothesis and the full conclusion? The supplied registry describes this as a metric-measure reconstruction theorem. The quoted passage conjectures an affirmative answer on the basis of numerical and heuristic arguments.

\textbf{Unresolved formulation.} The supplied record does not reproduce Theorem 3, its definitions, its other hypotheses, or its reconstruction conclusion, and does not identify the relevant contribution or theorem page within the report. Consequently, the underlying objects, domains, boundary conditions, quantifiers, and meaning of reconstruction cannot be specified from this evidence. The question above is a reference-dependent formulation, not a complete standalone mathematical conjecture.
\subsection{Short English statement}
The question asks whether the cited reconstruction theorem remains valid after its contractibility assumption is removed, with all other hypotheses and the conclusion retained. The source passage conjectures that it does, citing numerical and heuristic arguments. Theorem 3 itself is missing from the supplied record, so a complete mathematical statement requires further source clarification.
\subsection{Sources}
\begin{itemize}
\item Local UnsolvedMath registry, supplied record with permanent ID \(30005261\), title \emph{Removing Contractibility from Metric Measure Reconstruction}, exact source problem code \texttt{catalogue entry}; collection: \emph{Oberwolfach Reports Open Problems}. The record supplies proposed year 2022 but no named proposer.
\item Supplied report citation: \emph{Heat Kernels, Stochastic Processes and Functional Inequalities} (2023), \emph{Oberwolfach Reports}, pp.~2811--2893, DOI \texttt{10.4171/owr/2022/49}. \url{https://doi.org/10.4171/owr/2022/49}. The supplied passage asks about removing contractibility from Theorem 3; the theorem and contribution-level attribution are not supplied.
\item Supplied literature source: \emph{Oberwolfach Report 49}, \url{https://ems.press/content/serial-article-files/46987}. Literature checked at: \texttt{2026-08-22}. This is supplied dated curation, not an independently verified assessment. It labels the problem open and states that removal of contractibility is presented as a conjectural extension, with no proof or counterexample given in the source. The supplied background explicitly calls for independent verification of time-sensitive status and cited claims.
\end{itemize}
\end{document}
